\section{Devices and custody}
\label{sec:devices-and-custody}

A key file with no container placement is its own device. Distinct
\flag{--share-file} keys count as distinct physical devices, including when
\code{minimum_physical_devices} is greater than 1. Presenting the same
key twice still counts once.

\code{keyquorum-device} manages the container directory. \code{keyquorum
device} registers a slot's public key in the org store and binds it to the
opened container. Several identities can live as passphrase-wrapped slots
in one directory (a mounted USB, or a stand-in). The directory is an index:
\file{device.kq} (\code{KQDV} v2) is signed by \file{device.skey}, and each
slot token (\file{*.kqst}, \code{KQST}) seals that same device id. Opening a
slot checks the token against the descriptor, so a rewritten descriptor id
cannot make one container count as two devices. Path, mount, and volume
label are \emph{not} the device identity --- \code{relocate} moves a slot's
keypair onto another container, and that is the authenticated change of
\code{device_id}.

\begin{lstlisting}
keyquorum-device init ./usb
keyquorum-device provision ./usb --label M.S
keyquorum-device list ./usb
keyquorum device register ./usb --slot M.S --type encryption
keyquorum device bind ./usb --slot M.S
keyquorum-device relocate --from ./usb --to ./usb2 --label M.S
\end{lstlisting}

\code{register} records the public key; a lone key file still uses
\code{register} and \flag{--share-file}. \code{bind} writes the placement
from the container this process opened, tying that key to the container's
\code{device_id}.

\subsection{Presenting a slot}

Commands that unwrap a key also take \flag{--slot container=label} beside
\flag{--share-file}. This includes \code{reconstruct}, \code{add},
\code{bind}, \code{revoke}, an \code{access quorum} unlock,
\code{bridge private import}, \code{sign}, \code{remove-member}, and
\code{relay pull}. \code{tree countersign} takes \flag{--device} and
\flag{--slot} instead.

\begin{lstlisting}
keyquorum reconstruct <key-id> --slot ./usb=M.S --slot ./usb2=M.A \
  --output master.pub
keyquorum --db alice.sqlite relay pull --import --slot ./usb=M.S
\end{lstlisting}

\subsection{Custody modes}

\code{split}, \code{tree}, and \code{access quorum --state 0} set
\flag{--custody} (\code{hardware} or \code{logical}) and
\flag{--minimum-physical-devices}.

\begin{itemize}
  \item \textbf{Hardware} custody requires each share Shamir actually
  consumes to sit on its own device id.
  \item \textbf{Logical} custody lets several slots on one container meet
  that Shamir threshold together; the container still counts as one device
  toward the minimum. Logical mode is for development and constrained
  hardware --- logical slots are not a hardware quorum.
\end{itemize}

Reconstruction searches threshold-sized subsets for one that meets
\code{minimum_physical_devices}. Unused extra shares are not counted.

\begin{lstlisting}
keyquorum split --label master --threshold 2 --custody logical \
  --minimum-physical-devices 2 \
  --leaf M.S=SoftwareDepartment.pub --leaf M.A=AccountingDepartment.pub \
  --source master.pub --generate-keys --register
keyquorum tree <key-id> --custody hardware --minimum-physical-devices 2
\end{lstlisting}

\subsection{Parent approval}

\code{unlock_approval} is \code{none} (the default) or \code{parent}.
Parent approval is opt-in and sits on top of both Shamir and the device
count: each leaf that contributed a share needs its direct parent's
signature.

\begin{lstlisting}
# The leaf's own key, and a registered *signing* key for its direct
# parent M.S -- the approval check verifies against that registration.
keyquorum generate --type encryption --label M.S.1 --register \
  --public-key-out m-s-1.pub > m-s-1.key
keyquorum generate --type signing --label M.S --register \
  --public-key-out M.S.sign.pub > M.S.sign.key

keyquorum access quorum --state 0 --source ./secret.txt \
  --encrypted-path ./secret.txt.kqenc \
  --leaf M.S.1=m-s-1.pub --threshold 1 --unlock-approval parent
keyquorum access quorum --state 1 --id <file-id> --share-file m-s-1.key \
  --approve "M.S.1=M.S.sign.key"
\end{lstlisting}

\code{--approve} takes either \code{leaf=signing-key-file} or
\code{leaf=container>slot}.
